% ==============================================================================
% BAB II. PERSIAPAN PROGRAM DESA CANTIK
% ==============================================================================

\chapter{PERSIAPAN PROGRAM DESA CANTIK}

\section{Rancangan Kegiatan Desa Cantik 2026}

Pembinaan Desa Cantik di Desa Popontolen, Kecamatan Tumpaan Kabupaten Minahasa Selatan direncanakan untuk dilaksanakan melalui beberapa tahapan seperti pada Tabel~\ref{tab:rencana_kerja} Rencana Kerja Desa Cinta Statistik di bawah ini.

\begin{table}[htbp]
\centering
\small
\caption{Rencana Kerja Desa Cinta Statistik}
\label{tab:rencana_kerja}
\vspace{0.2cm}
\begin{tabularx}{\textwidth}{|c|Y|c|Y|}
\hline
\textbf{No} & \centering\textbf{Rincian Kegiatan} & \textbf{Tanggal Target} & \centering\textbf{Pelaksana} \tabularnewline
\hline
\textit{(1)} & \centering\textit{(2)} & \textit{(3)} & \centering\textit{(4)} \tabularnewline
\hline
1. & Koordinasi dengan Pemerintah Daerah & 16 April 2026 & Tim Desa Cantik BPS \tabularnewline
\hline
2. & Koordinasi dengan Pemerintah Kecamatan Tumpaan & 23 April 2026 & Tim Desa Cantik BPS \tabularnewline
\hline
3. & Koordinasi dengan Pemerintah Desa Popontolen & 23 April 2026 & Tim Desa Cantik BPS \tabularnewline
\hline
4. & Pemilihan Agen Desa Cantik Desa Popontolen & 06 Mei 2026 & Tim Desa Cantik BPS dan Pemerintah Desa Popontolen \tabularnewline
\hline
5. & Sosialisasi Desa Cantik Dengan Pemerintah Daerah & 07 Mei 2026 & Tim Desa Cantik BPS \tabularnewline
\hline
6. & Pencanangan Desa Cinta Statistik Desa Popontolen & 07 Mei 2026 & Tim Desa Cantik BPS \tabularnewline
\hline
7. & Identifikasi Desa Cinta Statistik Desa Popontolen & 07 Mei 2026 & Agen Desa Cantik Desa Popontolen \tabularnewline
\hline
8. & Rapat Bersama Agen Desa Popontolen dalam pembuatan Publikasi Infografis Desa dan pembuatan Media Sosial (FB atau IG) & 17 Juli 2026 & Agen Desa Cantik Desa Popontolen dan Tim Desa Cantik BPS Minsel \tabularnewline
\hline
9. & Rapat Agen Bersama Aparatur Desa & 30 Juli 2026 & Agen Desa Cantik Desa Popontolen \tabularnewline
\hline
\end{tabularx}
\end{table}

\section{Kegiatan Sosialisasi/Pencanangan}

Kegiatan Sosialisasi dan Pencanangan Desa Cinta Statistik (Desa Cantik) Desa Popontolen dilaksanakan pada tanggal \textbf{07 Mei 2026} di mana kegiatan dihadiri oleh Bupati Kabupaten Minahasa Selatan, Kepala Dinas PMD, Hukum Tua, Sekdes dan Kepala Jaga serta Agen Desa Cantik Desa Popontolen.

Kemudian kegiatan pencanangan Desa Cinta Statistik Desa Popontolen dicanangkan secara resmi oleh Kepala Dinas dan disaksikan oleh Ketua Tim Desa Cantik BPS Provinsi Sulawesi Utara, Kepala BPS Kabupaten Minahasa Selatan, Camat Tumpaan, dan Kepala Desa Popontolen beserta agen statistik dan perangkat desa (undangan, laporan dan dokumen lainnya terlampir).

\section{Kegiatan Pelatihan Pembina Statistik Desa Popontolen}

Kegiatan Pembinaan Statistik Desa di Desa Popontolen dilaksanakan dalam \textbf{2 (dua) tahap pembinaan utama} sebagai berikut:

\begin{enumerate}[label=\textbf{\arabic*.}]
  \item \textbf{Pembinaan Tahap I}
  \begin{itemize}
    \item \textbf{Hari / Tanggal:} Kamis, 7 Mei 2026
    \item \textbf{Tempat:} Kantor Desa Popontolen
    \item \textbf{Peserta:} Aparatur Desa, Agen Statistik Desa, serta Tokoh Masyarakat Desa Popontolen
    \item \textbf{Materi Pembinaan:}
    \begin{enumerate}[label=\arabic*.]
      \item \textbf{Identifikasi Kebutuhan Desa:} Pemetaan ketersediaan data, penentuan prioritas indikator statistik desa, dan identifikasi masalah statistik sektoral tingkat desa.
      \item \textbf{Pengumpulan Data:} Metodologi pengumpulan data sekunder maupun primer secara terstruktur di tingkat lingkungan/jaga.
      \item \textbf{Sosialisasi Kegiatan Sensus Ekonomi 2026:} Edukasi dan pemahaman awal mengenai urgensi, alur, serta persiapan pelaksanaan Sensus Ekonomi tahun 2026.
      \item \textbf{Sosialisasi DTSEN:} Pengenalan dan pentingnya Data Tunggal Sosial Ekonomi Nasional (DTSEN) untuk integrasi program perlindungan sosial dan pembangunan.
    \end{enumerate}
  \end{itemize}

  \item \textbf{Pembinaan Tahap II}
  \begin{itemize}
    \item \textbf{Hari / Tanggal:} Jumat, 17 Juli 2026
    \item \textbf{Tempat:} Kantor Desa Popontolen
    \item \textbf{Peserta:} Agen Statistik Desa dan Aparatur Desa Popontolen
    \item \textbf{Materi Pembinaan:}
    \begin{enumerate}[label=\arabic*.]
      \item \textbf{Cara Pengolahan Data:} Teknik pembersihan (\textit{data cleaning}), pengkodean, entri data, serta tabulasi data statistik desa menggunakan perangkat lunak/spreadsheet.
      \item \textbf{Cara Penyajian Data:} Penyusunan teknik visualisasi data yang informatif melalui pembuatan Monografi Desa, Publikasi Desa Dalam Angka, hingga penyusunan Infografis yang mudah dipahami oleh masyarakat umum.
    \end{enumerate}
  \end{itemize}
\end{enumerate}

\section{Anggaran Program Desa Cantik}

Kegiatan Desa Cinta Statistik (Desa Cantik) untuk anggaran masih menggunakan \textbf{DIPA BPS Kabupaten Minahasa Selatan}, di mana untuk ketersediaan anggaran di Desa belum ada anggaran mengenai kegiatan Desa Cantik (dokumen anggaran terlampir).
